\documentclass[11pt, a4paper]{article}
\usepackage[a4paper, margin=2.5cm]{geometry}
\usepackage{fontspec}
\usepackage[portuguese, english, bidi=basic, provide=*]{babel}
\babelprovide[import, onchar=ids fonts]{portuguese}
\babelfont{rm}{Noto Sans}
\usepackage{enumitem}
\usepackage{xcolor}
\usepackage{listings}
\usepackage[hidelinks]{hyperref}

% Definicao de cores para blocos de codigo e destaques
\definecolor{codebg}{RGB}{245, 247, 250}
\definecolor{codeframe}{RGB}{210, 215, 225}
\definecolor{primary}{RGB}{15, 76, 129}

% Configuracao do pacote listings para comandos e scripts
\lstset{
backgroundcolor=\color{codebg},
frame=single,
rulecolor=\color{codeframe},
basicstyle=\ttfamily\small,
breaklines=true,
columns=fullflexible,
keepspaces=true,
showstringspaces=false,
commentstyle=\color{gray}
}

\setlength{\parindent}{0pt}
\setlength{\parskip}{8pt}

\begin{document}

\begin{center}
{\LARGE \textbf{\color{primary} Guia do Professor: Preparando a ``AWS'' da Sala}} \\[8pt]
\rule{\textwidth}{1pt}
\end{center}

\vspace{10pt}

\section\*{1. Preparar o ambiente (antes dos alunos chegarem)}

No seu terminal Ubuntu, instale o Flask e o Docker (se já não tiver):

\begin{lstlisting}[language=bash]
pip install flask
\end{lstlisting}

\section\*{2. Ajustar o código}

Como você já criou o arquivo \texttt{material_de_apoio/codigo/criar_portal.py}, abra ele e certifique-se de que a última linha está assim:

\begin{lstlisting}[language=python]
if **name** == '**main**':
app.run(host='0.0.0.0', port=1961)
\end{lstlisting}

\section\*{3. Liberar a porta no Firewall (se estiver ativo)}

Se o seu Ubuntu tiver o firewall ativado (\texttt{ufw}), libere a porta para os alunos acessarem:

\begin{lstlisting}[language=bash]
sudo ufw allow 1961/tcp
\end{lstlisting}

\section\*{4. Criar a Rede Local (O ``Cabo da AWS'')}

Crie o hotspot Wi-Fi para os alunos se conectarem:

\begin{lstlisting}[language=bash]

# Substitua wlp2s0 pelo nome da sua placa Wi-Fi

sudo nmcli device wifi hotspot ifname wlp2s0 ssid "K8s-Cloud" password "SenhaSegura123"
\end{lstlisting}

Pegue o IP da sua máquina nessa rede:

\begin{lstlisting}[language=bash]
ip a show wlp2s0
\end{lstlisting}

\textit{(Suponha que o IP seja \texttt{10.42.0.1}. Escreva isso no quadro!)}

\section\*{5. Subir o Portal (em background)}

Para não travar o seu terminal (e você poder usar outros comandos durante a aula), rode o script em background usando o \texttt{nohup}:

\begin{lstlisting}[language=bash]
sudo nohup python3 material_de_apoio/codigo/criar_portal.py > portal.log 2>&1 &
\end{lstlisting}

Se quiser ver os logs de quem está se conectando, use: \texttt{tail -f portal.log}.

\section\*{6. Iniciar o Docker}

Certifique-se de que o serviço do Docker está rodando no seu PC:

\begin{lstlisting}[language=bash]
sudo systemctl start docker
\end{lstlisting}

\end{document}
